\chapter{Regime-Switching Models}
\label{cap:regime}
% ── index ───────────────────────────────────────────────────────
\index{Markov Switching|textbf}
\index{regime switching|textbf}
\index{markov@\texttt{markov()}|textbf}
\index{Hamilton, model of|textbf}
\index{transition probability|textbf}
\index{Hamilton filter}
\index{expected regime duration}

% ─────────────────────────────────────────────────────────────────────────────
\section{Structural heterogeneity over time}

Many economic series exhibit structural changes: different means,
volatilities or dynamics across distinct periods. Two approaches capture
this heterogeneity:

\begin{itemize}
  \item \textbf{Markov Switching} (Hamilton 1989): latent regime follows a
  Markov chain. The transition probability between regimes is estimated
  by the EM algorithm together with the parameters of each regime.
  \item \textbf{Threshold Regression} (Hansen 1997): regime change is
  determined by an observable variable (threshold variable $q$) crossing a
  threshold $\gamma$ to be estimated.
\end{itemize}

\subsection*{Markov-Switching AR(p)}

\[
  y_t = \mu_{S_t} + \sum_{k=1}^p \phi_k^{(S_t)}\,(y_{t-k} - \mu_{S_t}) + \sigma_{S_t}\,\varepsilon_t
\]

where $S_t \in \{1,\ldots,K\}$ is the latent regime with transition
probabilities $P(S_t = j \mid S_{t-1} = i) = p_{ij}$. The parameter $\mu_{S_t}$
is the regime intercept, $\phi_k^{(S_t)}$ are the AR coefficients and
$\sigma_{S_t}$ is the regime standard deviation.

% ─────────────────────────────────────────────────────────────────────────────
\section{DGP Verification}

DGP with two alternating regimes:
\begin{itemize}
  \item Regime 1 (expansion): $\mu_1 = 1{,}0$, $\phi_1 = 0{,}30$, $\sigma_1 = 0{,}5$
  \item Regime 2 (recession): $\mu_2 = -1{,}0$, $\phi_2 = 0{,}70$, $\sigma_2 = 1{,}5$
\end{itemize}

\begin{lstlisting}[caption={DGP Markov Switching AR(1) and estimation}]
seed(42)
generate df t      = _n
generate df regime = (t % 60 < 40)   // ~2/3 in regime 1
// y recursive with parameters conditional on regime
...
let m_ms = markov(df, y, k=2, p=1)
display m_ms
\end{lstlisting}

\subsection*{Markov Switching}

\begin{lstlisting}[language={}, basicstyle=\ttfamily\small,
  frame=none, numbers=none, backgroundcolor=\color{codbg},
  xleftmargin=0pt, xrightmargin=0pt, breaklines=false]
===================== Markov Switching AR(1) — 2 regimes =====================
Observations: 299   Log-L: -155.76   AIC: 327.51   BIC: 357.12
--------------------------- Transition Matrix --------------------------------
[ 0.9493   0.0507 ]
[ 0.0251   0.9749 ]
--------------------------- Regime 0 (Recession) -----------------------------
Intercept: -1.2864   AR.L1: 0.6169   Variance: 0.5141
--------------------------- Regime 1 (Expansion) -----------------------------
Intercept:  0.9946   AR.L1: 0.2956   Variance: 0.0625
--------------------------- Expected Durations --------------------------------
Regime 0: 19.7 periods    Regime 1: 39.9 periods
\end{lstlisting}

The model recovers the two regimes with good accuracy:
\begin{itemize}
  \item Regime 0 (recession): $\hat\mu = -1{,}29$ (true: $-1{,}0$),
  $\hat\phi = 0{,}62$ (true: 0{,}7). Expected duration 19{,}7 periods.
  \item Regime 1 (expansion): $\hat\mu = 0{,}99 \approx 1{,}0$,
  $\hat\phi = 0{,}30$ (true: 0{,}3). Expected duration 39{,}9 periods.
\end{itemize}

The transition matrix indicates high persistence of both regimes
($p_{00} = 0{,}95$, $p_{11} = 0{,}97$) — the series rarely switches regime,
consistent with the DGP ($60\%$ of the time in regime 1).

% ─────────────────────────────────────────────────────────────────────────────
\section{Syntax}

\begin{lstlisting}
markov(df, y, k=2, p=1)    // k regimes, AR of order p
\end{lstlisting}

\begin{center}
\begin{tabular}{lll}
\toprule
\textbf{Option} & \textbf{Default} & \textbf{Description} \\
\midrule
\texttt{k} & 2 & number of regimes \\
\texttt{p} & 1 & AR order within each regime \\
\bottomrule
\end{tabular}
\end{center}

\begin{nota}
  The number of regimes $k$ is fixed and must be specified a priori.
  Information criteria (AIC, BIC) compared across models with $k=2$ and
  $k=3$ guide this choice. For data with $N < 200$, $k > 3$ is rarely
  identified reliably.
\end{nota}
